% Target: rmp-workbench-iii-mpo-on-peps-definition
% Formula: The source illustrates a PEPS tensor and its four-site MPO projector.
% Ink: Kernel flat frame; a five-leg PEPS star and four boxes on a red MPO ring.
% Boundary: The PEPS star has five open legs; the MPO has eight virtual openings.
%   The source equality is schematic, not a boundary-compatible equation claim.
% Source: author source ImagesReview Section III/ImagesReview Section III.tex
%   lines 521-537; archival output DefofMPOonjPEPSten.pdf
% Capabilities: kernel, typed-ports
\[
  \tenkzkernel
  \tndeclare{species}{operator}{hue=source:red}
  \begin{tenkz}[rows={wire}, cols=1, bonds=none]
    \tn[skin=dot, ports={0:virtual, 180:virtual,
        45:virtual, 225:virtual, 90:physical}]{}
  \end{tenkz}
  \;=\;
  \begin{tenkz}[rows={wire,wire,wire}, cols=3, bonds=none]
    \tn[skin=mpo, at=(2,1), name=w,
        ports={0:virtual, 90:virtual, 180:virtual, 270:virtual}]{}
    \tn[skin=mpo, at=(1,2), name=n,
        ports={0:virtual, 90:virtual, 180:virtual, 270:virtual}]{}
    \tn[skin=mpo, at=(2,3), name=e,
        ports={0:virtual, 90:virtual, 180:virtual, 270:virtual}]{}
    \tn[skin=mpo, at=(3,2), name=s,
        ports={0:virtual, 90:virtual, 180:virtual, 270:virtual}]{}
    \tnwire[route=arc, species=operator]{w.90}{n.180}
    \tnwire[route=arc, species=operator]{n.0}{e.90}
    \tnwire[route=arc, species=operator]{e.270}{s.0}
    \tnwire[route=arc, species=operator]{s.180}{w.270}
  \end{tenkz}
\]
